\documentclass[conference]{IEEEtran}
\IEEEoverridecommandlockouts
% The preceding line is only needed to identify funding in the first footnote. If that is unneeded, please comment it out.
\usepackage{cite}
\usepackage{amsmath,amssymb,amsfonts}
\usepackage{graphicx}
\usepackage{textcomp}
\usepackage{xcolor}
\usepackage{hyperref}
\def\BibTeX{{\rm B\kern-.05em{\sc i\kern-.025em b}\kern-.08em
    T\kern-.1667em\lower.7ex\hbox{E}\kern-.125emX}}
\begin{document}

\title{\textbf{FoodHub} \\ \small{Real-Time Food Delivery System}}

\author{\IEEEauthorblockN{Student Name: Dhruv Dhangar}
\IEEEauthorblockA{Enrollment ID: \textbf{202201144} \\
BTP Mode: \textbf{\textit{Off Campus / ITP}} \\
\textit{School of Technology, Dhirubhai Ambani University, }\\
Gandhinagar, India \\
202201144@dau.ac.in}
\and
\IEEEauthorblockN{Mentor's Name: Varsha Oberoi}
\IEEEauthorblockA{Company's Name: \textit{Radixweb} \\
Company's Address: Ekyarth, Malabar County Rd \\
Chharodi, Ahmedabad - 382481, India \\
varsha.oberoi@radixweb.com \\
On-Campus Mentor: BTP Coordinator}
}

\maketitle

\begin{abstract}
This report details the development of FoodHub, a real-time food delivery web application. The primary goal of the project was to build a functional platform where users can browse local restaurants, add items to a digital cart, and place orders securely. The application uses a decoupled client-server architecture. The frontend is built with Vite and TypeScript \cite{b2}, while the backend relies on Node.js \cite{b1} and Express.js \cite{b4}. For data storage, we chose MongoDB \cite{b3} to handle the complex, nested menu and order data flexibly. We also implemented JSON Web Tokens (JWT) \cite{b7} for secure user authentication. Throughout this report, we discuss the overall system design, the specific algorithms used (such as fuzzy search for the menu and the Luhn algorithm for checkout validation), and the practical lessons learned while building and deploying a full-stack project from scratch.
\end{abstract}

\begin{IEEEkeywords}
Food Delivery, Full-Stack Development, TypeScript, Node.js, Express.js, MongoDB, Web Application, Authentication, API Design.
\end{IEEEkeywords}

\section{Introduction}
Online food delivery has become a standard service for the restaurant industry today. The goal of this project was to build a working food delivery application from the ground up to understand how these systems operate behind the scenes. Instead of building a single monolithic application where everything is tangled together, we decided to use a modern client-server architecture, cleanly splitting the frontend and backend.

The frontend uses Vite to bundle our TypeScript code \cite{b2}. This setup makes the development process faster and helps catch errors early through strict type checking. The backend is a separate Node.js server using the Express framework \cite{b4}, which is responsible for managing the core business logic, processing incoming orders, and interacting with the database.

By breaking the system into these two distinct parts, we ensure the application is much easier to maintain and debug. It also makes it easier to scale the application if more users start placing orders simultaneously. The final system includes essential features like dynamic data fetching, real-time order tracking, and secure user logins using JSON Web Tokens (JWT) \cite{b7}.

\section{System Architecture}
We designed the FoodHub application using a decoupled approach. This means the user interface operates entirely independently from the server-side logic, communicating only via API requests.

\subsection{Client-Side Architecture (Frontend)}
The frontend is built as a Single Page Application (SPA). This approach prevents the browser from reloading the entire page every time the user clicks a link, which makes the app feel much faster and more responsive.

\begin{itemize}
\item \textbf{Technology Stack:} We used TypeScript \cite{b2}, plain HTML5, and standard CSS, along with Bootstrap 5 to help with layout and responsive design.
\item \textbf{Design Pattern:} The code is organized into modular service classes (like \textit{MenuService} and \textit{CartService}). These classes handle fetching data from the API so that the user interface code doesn't get cluttered with network requests.
\item \textbf{User Interface:} The app includes a home page for restaurant discovery, a menu page with category filters, a shopping cart that saves your items using local storage, a profile section for order history, and a secure checkout page.
\end{itemize}

\subsection{Server-Side Architecture (Backend)}
The backend acts as the brain of the application. It processes incoming requests, checks if they are valid, and reads or writes data to the database.

\begin{itemize}
\item \textbf{Technology Stack:} The server is built with Node.js \cite{b1} and Express.js \cite{b4}. We also wrote the backend in TypeScript to keep our data types consistent across both sides of the project.
\item \textbf{Routing and Controllers:} We split the API endpoints into different routers. For example, all order-related requests go to the \textit{orderController}. This keeps the codebase organized and easier to read.
\item \textbf{Security:} We added middleware to handle Cross-Origin Resource Sharing (CORS), validate incoming data, and hash user passwords using the \textit{bcryptjs} library before saving them. We also set up custom middleware to verify JWTs for routes that require the user to be logged in \cite{b7}.
\end{itemize}

\subsection{Database Schema}
We used MongoDB \cite{b3}, which is a NoSQL database. It is a great fit for an e-commerce app because things like restaurant menus can be stored as nested objects rather than rigid, flat tables. We used Mongoose to interact with MongoDB and enforce rules on our data.

\begin{itemize}
\item \textbf{Users Collection:} This stores user details and hashed passwords. We made sure email addresses are stored as unique, case-insensitive values so users can't accidentally create duplicate accounts.
\item \textbf{Restaurants and Menus Collection:} This stores restaurant details and their food items. Because MongoDB allows nested data, we can pull a restaurant and its entire menu in a single, fast database request.
\item \textbf{Orders Collection:} This tracks user purchases. It saves the items bought, the total cost, the delivery address, and the current status of the order (like ``Preparing'' or ``Delivered'').
\end{itemize}

\section{Key Methodologies and Algorithms}
While building the app, we implemented several specific algorithms to improve how the system handles user input and security.

\subsection{Fuzzy Search Algorithm}
We wrote a custom fuzzy search function for the restaurant menu. Instead of just looking for exact word matches, this algorithm checks how similar the user's search string is to the food item names. This means that if a user makes a typo while searching, the app can still figure out what they are looking for and show the right results.

\subsection{Luhn Algorithm for Credit Card Validation}
To make the checkout process smoother, we implemented the Luhn algorithm \cite{b6} on the frontend. This is a mathematical formula used to check if a credit card number is structurally valid. By running this check directly in the user's browser, we can immediately tell them if they made a typo, without needing to send bad data to the server first.

\subsection{JWT Authentication and Session Management}
For user logins, we decided to use JSON Web Tokens \cite{b7}. When a user logs in successfully, the backend creates a token and sends it back to the browser. If the user selects ``Remember Me,'' we save this token in \textit{localStorage} so they don't have to log in again later. If they don't, we put it in \textit{sessionStorage} so it deletes automatically when they close the tab. Every time the app needs to fetch private data, it simply attaches this token to the request header.

\subsection{State Persistence and Real-Time Updates}
One of the challenges of order tracking is making sure the status doesn't reset if the user refreshes the page. We solved this by having the frontend periodically check and sync the order status with the backend. This way, the server always holds the true state of the order, and the frontend simply displays what the server tells it.

\section{Deployment and Infrastructure}
Once the code was finished, we deployed the application to live cloud servers so it can be accessed over the internet.

\begin{itemize}
\item \textbf{Frontend Hosting (Vercel):} We deployed the frontend application on Vercel \cite{b5}. Vercel makes it very easy to host frontend projects, and it automatically updates the live website whenever we push new code to our GitHub repository.
\item \textbf{Backend Hosting (Render):} We hosted the Node.js server on Render. Render gives us a safe place to store our secret environment variables (like the MongoDB connection string and the JWT secret key) and keeps the server running continuously.
\end{itemize}

\section{Learning Outcomes}
Building this project from scratch was a huge learning experience. It required understanding how all the different pieces of a modern web application fit together.

\subsection{Technical Skills Acquired}
\begin{itemize}
\item \textbf{Full-Stack Integration:} We learned how to properly connect a frontend interface to a backend API and how to fix common networking issues like CORS errors.
\item \textbf{Asynchronous Programming:} We got a lot of practice using JavaScript Promises and \texttt{async/await} to handle network requests without freezing the web page.
\item \textbf{Database Management:} We learned how to design a NoSQL database schema and write practical queries using Mongoose \cite{b3}.
\item \textbf{Deployment:} We figured out how to take code from a local development environment and put it on live servers like Vercel and Render.
\end{itemize}

\subsection{Non-Technical Skills Acquired}
\begin{itemize}
\item \textbf{Planning:} We learned that designing the database structure and planning the API endpoints on paper before writing code saves a massive amount of time later.
\item \textbf{Debugging:} We significantly improved our ability to track down bugs, especially when data wasn't passing correctly between the frontend and the backend.
\item \textbf{Project Lifecycle:} We experienced the full software development lifecycle, from the initial idea all the way to a live production deployment.
\end{itemize}

\section{Contributions}
Throughout the development process, several major tasks were completed to bring the project to life:

\begin{itemize}
\item \textbf{UI/UX Engineering:} We designed and built a responsive, custom user interface.
\item \textbf{Frontend Architecture:} We set up the TypeScript project \cite{b2} and organized the code into reusable service classes.
\item \textbf{Backend Development:} We wrote the Express server \cite{b4}, creating all the necessary API routes, models, and controllers.
\item \textbf{Security:} We implemented password hashing and the complete JWT login system \cite{b7}.
\item \textbf{Algorithmic Work:} We wrote the logic for the fuzzy menu search and the client-side credit card validation \cite{b6}.
\item \textbf{Database Setup:} We wrote scripts to automatically fill the database with sample restaurants and menu items so we could test the app thoroughly.
\item \textbf{Deployment:} We successfully deployed both parts of the application to the internet.
\end{itemize}

\section{Conclusion and Future Scope}
The FoodHub project successfully demonstrates how a modern, full-stack web application is built. We managed to create a system that handles user accounts, displays dynamic data, processes simulated orders, and tracks deliveries. Breaking the project into a separate frontend and backend proved to be a great decision, as it kept development highly organized.

Looking forward, there are a few features we would like to add:
\begin{itemize}
\item \textbf{WebSockets:} Right now, the frontend asks the server for order updates by polling. We'd like to replace this with WebSockets so the server can push updates to the browser instantly.
\item \textbf{Real Payments:} We currently simulate the checkout process. In the future, we want to integrate a live payment provider like Stripe or Razorpay.
\item \textbf{Recommendations:} It would be interesting to use basic machine learning to suggest food items to users based on their past purchase history.
\end{itemize}

\begin{thebibliography}{00}
\bibitem{b1} Node.js Foundation, ``Node.js Documentation,'' [Online]. Available: https://nodejs.org/en/docs/.
\bibitem{b2} Microsoft, ``TypeScript Handbook,'' [Online]. Available: https://www.typescriptlang.org/docs/handbook/intro.html.
\bibitem{b3} MongoDB Inc., ``MongoDB Manual,'' [Online]. Available: https://www.mongodb.com/docs/manual/.
\bibitem{b4} Express.js, ``Express - Node.js web application framework,'' [Online]. Available: https://expressjs.com/.
\bibitem{b5} Vercel, ``Vercel Documentation,'' [Online]. Available: https://vercel.com/docs.
\bibitem{b6} H. P. Luhn, ``Computer for Verifying Numbers,'' U.S. Patent 2,950,048, Aug. 23, 1960.
\bibitem{b7} M. Jones, J. Bradley, and N. Sakimura, ``JSON Web Token (JWT),'' RFC 7519, May 2015.
\end{thebibliography}

\end{document}
